% Avambraccio orizzontale che regge un peso: leva di terzo genere con il fulcro
% nel gomito. Il bicipite tira a b1 dal gomito, il peso agisce a b2.
\begin{tikzpicture}[>=Stealth, font=\small]
  % omero (verticale) e avambraccio (orizzontale)
  \filldraw[fill=esempio!15, draw=esempio, thick, rounded corners=3pt] (-0.25,0) rectangle (0.25,3.2);
  \filldraw[fill=esempio!25, draw=esempio, thick, rounded corners=3pt] (-0.2,-0.22) rectangle (6.2,0.22);
  \fill[richiamo!70] (0,0) circle (0.2);
  \fill (0,0) circle (1.5pt);
  \node[below left] at (-0.2,-0.2) {gomito $O$};
  % bicipite
  \draw[medica!35, line width=5pt, line cap=round] (0.3,2.9) .. controls (1.0,2.2) and (1.1,1.2) .. (1.0,0.22);
  \node[medica, right] at (1.1,2.3) {bicipite};
  \draw[->, medica!80!black, very thick] (1.0,0.22) -- (1.0,1.5) node[right] {$\vec F_b$};
  % peso nella mano
  \draw[thick] (6.0,-0.22) -- (6.0,-0.6);
  \filldraw[fill=richiamo!60, draw=black] (5.7,-0.6) rectangle (6.3,-1.05);
  \draw[->, attenzione, very thick] (6.0,-1.05) -- (6.0,-2.0) node[right] {$\vec F_p$};
  % distanze
  % peso dell'avambraccio nel suo baricentro
  \draw[->, attenzione!70, thick] (2.6,-0.22) -- (2.6,-1.0) node[below] {$\vec F_{av}$};
  % distanze
  \draw[dashed, richiamo] (0,-0.25) -- (0,-2.55);
  \draw[|<->|, formula] (0,-0.6) -- (1.0,-0.6) node[midway, below] {$b_1$};
  \draw[|<->|, formula] (0,-2.4) -- (6.0,-2.4) node[midway, below] {$b_2$};
\end{tikzpicture}
